\chapter{Całka McShane'a i przestrzenie funkcji BV oraz AC}
\providecommand{\atlasref}[2]{#2}
\label{ma1:ch:calka-mcshane-a-i-przestrzenie-bv-ac}
\begingroup
\providecommand{\osc}{\operatorname{osc}}
\providecommand{\Var}{\operatorname{Var}}
\providecommand{\mesh}{\operatorname{mesh}}
\providecommand{\supp}{\operatorname{supp}}
\providecommand{\AC}{\mathrm{AC}}
\providecommand{\BV}{\mathrm{BV}}
\providecommand{\HK}{\mathrm{KH}}
\providecommand{\ind}[1]{\chi_{#1}}
\providecommand{\Dn}{\underline D}
\providecommand{\Du}{\overline D}
\providecommand{\KS}{\mathrm{KS}}
Przyjmujemy oznaczenia i konwencje z rozdziału \ref{ma1:ch:kurzweil-henstock}.
\section{Całki McShane'a i Lebesgue'a}
\begin{definicja}\label{it:def:mcshane}
Niech $a<b$ i niech $f:[a,b]\to\R$ ma skończone wartości. Dla podziału $a=x_0<\dots<x_n=b$ rodzinę $\mathcal P=\{(t_i,[x_{i-1},x_i])\}_{i=1}^n$ nazywamy podziałem znaczonym; jest on jest \emph{wolny $\delta$-drobny}, gdy $t_i\in[a,b]$ oraz $[x_{i-1},x_i]\subset(t_i-\delta(t_i),t_i+\delta(t_i))$, bez wymagania $t_i\in[x_{i-1},x_i]$. Funkcja $f$ jest całkowalna w sensie McShane'a, $f\in\mathcal M[a,b]$, jeśli istnieje $I\in\R$ takie, że dla każdego $\varepsilon>0$ istnieje cechowanie $\delta:[a,b]\to(0,\infty)$, dla którego każdy wolny $\delta$-drobny podział spełnia
\[
\left|S(f,\mathcal P)-I\right|<\varepsilon,\qquad
 S(f,\mathcal P)=\sum_i f(t_i)(x_i-x_{i-1}).
\]
Piszemy wówczas $I=(\mathcal M)\int_a^b f$. Na przedziale jednopunktowym przyjmujemy całkę zero.
\end{definicja}
\begin{stwierdzenie}[Bezpośrednie porównanie całek Riemanna i McShane'a]\label{ma1:kh:prop:riemann-mcshane}
$\mathcal R[a,b]\subset\mathcal M[a,b]\subset\HK[a,b]$, a wartości całek są zgodne.
\end{stwierdzenie}
\begin{proof}
Druga inkluzja wynika z warunków na znaczniki. Dla $a=b$ przyjmujemy całki równe zeru. Dalej $a<b$. Dla pierwszej niech $f\in\mathcal R[a,b]$ i $\varepsilon>0$. Ustalmy $M=1+\|f\|_\infty$ oraz, na mocy równoważności całkowalności Riemanna i Darboux
(twierdzenie \ref{ma1:thm:riemann-darboux-equivalence}), podział $P_0$ z $U(f,P_0)-L(f,P_0)<\varepsilon/2$. Niech $B$ będzie skończonym zbiorem punktów jego podziału. Dla $t\notin B$ wybieramy $\delta(t)<\operatorname{dist}(t,B)$, zaś w punktach $t\in B$ dobieramy dodatnie cechowania tak małe, aby $4M\sum_{t\in B}\delta(t)<\varepsilon/2$. Każdy przedział wolnego podziału ze znacznikiem poza $B$ leży w tej samej komórce $P_0$ co jego znacznik. Oznaczmy komórki $P_0$ przez $J_k$, a oscylację $f$ na $J_k$ przez $\omega_k=\sup_{J_k}f-\inf_{J_k}f$. Z lokalności całki Riemanna (twierdzenie \ref{ma1:thm:riemann-additivity}) mamy dla każdej komórki $I_i$
\[
 S(f,\mathcal P)-\int_a^b f=\sum_i\left(f(t_i)|I_i|-\int_{I_i} f\right).
\]
Gdy $t_i\notin B$ oraz $t_i\in J_k$, $I_i\subset J_k$, oszacowanie całki przez kresy funkcji daje moduł składnika nie większy niż $\omega_k|I_i|$. Łączna długość tych komórek w $J_k$ jest nie większa niż $|J_k|$, więc suma błędów dla znaczników spoza $B$ jest nie większa niż $\sum_k\omega_k|J_k|=U(f,P_0)-L(f,P_0)$. Dla ustalonego $t\in B$ wszystkie komórki z tym znacznikiem mają rozłączne wnętrza i leżą w otoczeniu długości $2\delta(t)$, nawet jeśli znacznik do nich nie należy. Moduł każdego składnika jest nie większy niż $2M|I_i|$, więc ich łączny błąd nie przekracza $4M\delta(t)$. Dodając oszacowania, dostajemy błąd mniejszy niż $\varepsilon$.
\end{proof}

\begin{twierdzenie}[Kryterium Cauchy'ego, liniowość i porządek]\label{it:thm:mcshane-basic}
Dla funkcji $f:[a,b]\to\R$ o skończonych wartościach
całka McShane'a istnieje wtedy i tylko wtedy, gdy dla każdego
$\varepsilon>0$ istnieje cechowanie $\delta$ takie, że
$|S(f,\mathcal P)-S(f,\mathcal Q)|<\varepsilon$ dla wszystkich
wolnych $\delta$-drobnych podziałów $\mathcal P,\mathcal Q$.
Całka jest jednoznaczna, liniowa i monotoniczna.
Na przedziale jednopunktowym przyjmujemy całkę zero.
\end{twierdzenie}
\begin{proof}
Dla $a=b$ oba rodzaje sum są puste i wszystkie tezy wynikają z konwencji całki zerowej. Dalej zakładamy $a<b$. Każdy podział KH jest wolnym podziałem, zatem lemat Cousina
\ref{ma1:kh:lem:kh-cousin} zapewnia istnienie wolnego podziału
dla dowolnego cechowania. Konieczność kryterium wynika
z nierówności trójkąta. Dla dostateczności wybieramy cechowania
$\delta_n$ dla błędów $2^{-n}$, zastępując je kolejno
minimami z poprzednimi. Wybieramy wolne $\delta_n$-drobne
podziały $\mathcal P_n$. Dla $m\ge n$ obie sumy są drobne
względem $\delta_n$, więc ich różnica ma moduł mniejszy
od $2^{-n}$. Sumy tworzą ciąg Cauchy'ego o granicy $I$.
Porównując dowolny wolny $\delta_n$-drobny podział
z $\mathcal P_m$ dla $m\ge n$ i przechodząc do granicy,
otrzymujemy błąd względem $I$ nie większy od $2^{-n}$. Dla dowolnej wymaganej dokładności $\varepsilon>0$ wybieramy $n$ takie, że $2^{-n}<\varepsilon$. Cechowanie $\delta_n$ spełnia wtedy definicję całki McShane'a z wartością $I$.

Jeśli dwie liczby spełniają definicję całki, wybieramy
wspólne cechowanie dla błędu $\varepsilon/2$ i jeden drobny
podział; różnica tych liczb ma moduł mniejszy od $\varepsilon$.
Dla dwóch funkcji całkowalnych $f,g$ i współczynników $\alpha,\beta\in\R$ wybieramy minimum ich cechowań dla błędu $\varepsilon/(1+|\alpha|+|\beta|)$. Liniowość skończonych sum daje
\[
 |S(\alpha f+\beta g,\mathcal P)-\alpha I_f-\beta I_g|
 \le |\alpha||S(f,\mathcal P)-I_f|
     +|\beta||S(g,\mathcal P)-I_g|<\varepsilon.
\]
To dowodzi istnienia całki kombinacji liniowej oraz jej wartości.
Dla $f\le g$ mamy $S(f,\mathcal P)\le S(g,\mathcal P)$
na każdym podziale, ponieważ długości komórek są nieujemne.
Dla $\varepsilon>0$ wybieramy wspólne cechowanie dla błędu $\varepsilon/2$ i wolny drobny podział (istniejący na mocy lematu Cousina). Wtedy $I_f\le S(f,\mathcal P)+\varepsilon/2\le S(g,\mathcal P)+\varepsilon/2\le I_g+\varepsilon$. Dowolność $\varepsilon$ daje $I_f\le I_g$.
\end{proof}

\begin{wniosek}[Zmiany na zbiorze zerowym dla całki McShane'a]
\label{ma1:kh:cor:mcshane-null-modification}
Jeżeli $h:[a,b]\to\R$ ma skończone wartości i jest zerowa
poza zbiorem miary Lebesgue'a zero, to jej całka McShane'a
istnieje i wynosi zero. Zmiana funkcji na zbiorze zerowym
nie zmienia jej całkowalności w sensie McShane'a ani wartości całki.
\end{wniosek}
\begin{proof}
Stosujemy cechowanie z dowodu lematu
\ref{ma1:kh:lem:kh-null-modification}.
Dla wolnego drobnego podziału, tak samo jak dla podziału KH,
komórki o znacznikach z $N_k$ mają rozłączne wnętrza
i leżą w otwartym zbiorze $U_k$.
Ich łączna długość nie przekracza $\lambda_1(U_k)$;
nie potrzebujemy przynależności znacznika do komórki.
Zatem moduł całej sumy jest ograniczony przez
$\sum_{k\ge1} k\lambda_1(U_k)<\varepsilon$.
Ostatnia teza wynika z liniowości
(twierdzenie \ref{it:thm:mcshane-basic}) zastosowanej do różnicy funkcji,
w obu kierunkach.
\end{proof}

\begin{twierdzenie}[Bezwzględna całkowalność]\label{it:thm:mcshane-absolute}
Jeżeli $f\in\mathcal M[a,b]$, to $|f|\in\mathcal M[a,b]$. Tak więc całka McShane'a jest bezwzględna, mimo że definiujemy ją przez sumy znaczonych przedziałów.
\end{twierdzenie}
\begin{proof}
Niech $I=\int_a^b f$ w sensie McShane'a. Wybieramy cechowanie $\delta$ zapewniające błąd mniejszy niż $\varepsilon/3$. Dla dwóch dowolnych wolnych $\delta$-drobnych podziałów $\mathcal P=\{(t_i,I_i)\}$ i $\mathcal Q=\{(s_j,J_j)\}$ rozważamy wspólne rozdrobnienie $K_{ij}=I_i\cap J_j$, pomijając przecięcia zerowej długości. Każda komórka $K_{ij}$ leży w otoczeniach cechowania zarówno $t_i$, jak i $s_j$, więc obydwa znaczniki są legalne. Wybierając w jednej sumie znacznik z większą wartością $f$, a w drugiej z mniejszą, otrzymujemy dwie dopuszczalne sumy $S_+,S_-$, zatem
\[
\sum_{i,j}|f(t_i)-f(s_j)|\,|K_{ij}|=S_+-S_-<2\varepsilon/3.
\]
Ponieważ $|I_i|=\sum_j|K_{ij}|$, $|J_j|=\sum_i|K_{ij}|$ i $\bigl||u|-|v|\bigr|\le|u-v|$, mamy
$|S(|f|,\mathcal P)-S(|f|,\mathcal Q)|\le\sum_{i,j}|f(t_i)-f(s_j)|\,|K_{ij}|<\varepsilon$.
Kryterium Cauchy'ego dla wolnych podziałów z twierdzenia \ref{it:thm:mcshane-basic} dowodzi istnienia całki McShane'a z $|f|$. Dla KH podobny argument nie działa: znacznik pierwotnej komórki nie musi należeć do jej przecięcia z inną komórką, co narusza warunek KH.
\end{proof}
\begin{lemat}[Bezwzględna ciągłość całki]\label{ma1:kh:lem:integral-ac}\label{it:lem:integral-absolute-continuity}
Jeśli $f\in\mathcal L^1(\mu)$, to dla każdego $\varepsilon>0$ istnieje $\eta>0$ takie, że $\mu(A)<\eta$ pociąga $\int_A|f|\,\mathrm d\mu<\varepsilon$ dla każdego $A\in\mathcal F$ w przestrzeni z miarą $(X,\mathcal F,\mu)$.
\end{lemat}
\begin{proof}
Jest to wniosek
\atlasref{ME-17::cor:uniform-abs-cont-int-17}{ME-17: ciągłość całki względem zbioru całkowania}.
Funkcja $f\in\mathcal L^1(\mu)$ wyznacza klasę w $L^1(\mu)$,
więc spełnia jego założenia. Nie wymaga się skończoności
ani $\sigma$-skończoności miary $\mu$.
\end{proof}
\begin{twierdzenie}[McShane = Lebesgue]\label{it:thm:mcshane-lebesgue}
Niech $a\le b$. Dla funkcji rzeczywistych zachodzi
$\mathcal M[a,b]=\mathcal L^1([a,b],\lambda_1)$
jako równość zbiorów konkretnych funkcji o skończonych
wartościach, gdzie $\lambda_1$ jest miarą uzupełnioną.
Dowolne zmiany skończonych wartości na zbiorach
miary zero zachowują całkowalność i wartość całki.
Wartości obu całek są równe.
\end{twierdzenie}
\begin{proof}
Dla $a=b$ każda skończona funkcja na singletonie
jest mierzalna, całkowalna i obie całki wynoszą zero.
Dalej zakładamy $a<b$.

\emph{Krok 1: Lebesgue $\Rightarrow$ McShane.} Niech $f\in\mathcal L^1$ i $\varepsilon>0$. Wybieramy $h>0$ z $h(b-a)<\varepsilon/3$ i rozbijamy dziedzinę na mierzalne zbiory $E_k=\{kh\le f<(k+1)h\}$, $k\in\mathbb Z$. Z lematu \ref{ma1:kh:lem:integral-ac} wybieramy $\eta>0$ tak, że $\lambda_1(A)<\eta$ pociąga $\int_A|f|\,\mathrm d\lambda_1<\varepsilon/3$. Ustalamy dodatnie liczby $w_k=2^{-|k|}/3$ ($k\in\Z$), których suma wynosi $1$. Każdy $E_k$ jest mierzalny w $\mathbb R$ i ma miarę co najwyżej $b-a$. Z \atlasref{ME-04::thm:lebesgue-outer-regular}{ME-04: aproksymacji otwartymi zbiorami} wybieramy otwarte $O_k\supset E_k$ w $\mathbb R$, a następnie bierzemy $G_k=O_k\cap[a,b]$. Możemy tak dobrać $O_k$, aby
\[
 \lambda_1(G_k\setminus E_k)<
 w_k\min\left\{\eta/2,\frac{\varepsilon}{6(1+|kh|+h)}\right\}.
\]
Z przeliczalnej podaddytywności miary i $\sum_{k\in\mathbb Z}w_k=1$ dla $U=\bigcup_k(G_k\setminus E_k)$ mamy wtedy $\lambda_1(U)\le\eta/2$ i $\int_U|f|\,\mathrm d\lambda_1<\varepsilon/3$, a suma $\sum_k(|kh|+h)\lambda_1(G_k\setminus E_k)$ nie przekracza $\varepsilon/6$. Dla $t\in E_k$ wybieramy $\delta(t)>0$ z $(t-\delta(t),t+\delta(t))\cap[a,b]\subset G_k$.

Niech $\mathcal P=\{(t_i,I_i)\}$ będzie dowolnym wolnym $\delta$-drobnym podziałem. Zatem $I_i\subset G_k$ dla $t_i\in E_k$, również gdy $t_i\notin I_i$. Ponieważ komórki mają rozłączne wnętrza, a ich wspólne końce są zbiorami zerowymi, z addytywności całki Lebesgue'a mamy
\[
 S(f,\mathcal P)-\int_a^b f(x)\,\mathrm d\lambda_1(x)
 =\sum_i\int_{I_i}\bigl(f(t_i)-f(x)\bigr)\,\mathrm d\lambda_1(x).
\]
Przy znaczniku $t_i\in E_k$ na części $I_i\cap E_k$ moduł różnicy wartości $f(t_i)$ i $f(x)$ jest mniejszy od $h$, ponieważ obie wartości należą do $[kh,(k+1)h)$. Pozostałe mierzalne części $I_i\setminus E_k$ leżą w $U$; ich przecięcia zawierają co najwyżej wspólne końce komórek, a więc mają miarę zero. Dla ustalonego $k$ suma ich miar po komórkach ze znacznikami w $E_k$ nie przekracza $\lambda_1(G_k\setminus E_k)$. Po rozdzieleniu sumy otrzymujemy
\[
\left|S(f,\mathcal P)-\int_a^b f\,\mathrm d\lambda_1\right|
\le h(b-a)+\sum_k(|kh|+h)\lambda_1(G_k\setminus E_k)+\int_U|f|\,\mathrm d\lambda_1<\varepsilon.
\]

\emph{Krok 2: McShane $\Rightarrow$ Lebesgue.} Ponieważ każdy podział ze znacznikami należącymi do komórek jest również podziałem z wolnymi znacznikami, całkowalność McShane'a pociąga całkowalność KH z tą samą wartością całki. Mierzalność $f$ wynika z twierdzenia \ref{ma1:kh:thm:kh-primitive-derivative}. Z twierdzenia \ref{it:thm:mcshane-absolute} także $|f|\in\mathcal M$, z całką $J\ge0$. Obcięcia $h_n=\min(|f|,n)$ są ograniczone i mierzalne, zatem Lebesgue'a całkowalne, a na mocy kroku 1 także McShane'a całkowalne. Monotoniczność całki McShane'a (twierdzenie \ref{it:thm:mcshane-basic}) i zgodność z kroku 1 dają $\int h_n\,\mathrm d\lambda_1\le J$. Ponieważ $0\le h_n\uparrow |f|$ punktowo, twierdzenie \atlasref{ME-08::thm:mct}{ME-08} daje $\int_{[a,b]}|f|\,\mathrm d\lambda_1=\lim_n\int_{[a,b]}h_n\,\mathrm d\lambda_1\le J<\infty$. Krok 1 i jednoznaczność całki McShane'a
(twierdzenie \ref{it:thm:mcshane-basic})
dowodzą zgodności wartości. Jeżeli zmienimy
skończone wartości $f$ na zbiorze miary zero,
to zupełność miary zapewnia mierzalność zmodyfikowanej
funkcji, a twierdzenie \atlasref{ME-07::thm:ae-integral-invariance}{ME-07: niezmienniczość całki przy równości prawie wszędzie}
zapewnia równość całek Lebesgue'a i zachowanie
całkowalności. Z wykazanej równoważności wynika to
samo dla McShane'a. Ponieważ każdy podział KH jest
wolnym podziałem, w szczególności dla każdej
$f\in\mathcal L^1$ zachodzi $f\in\HK$ i
$(\HK)\int_a^b f=\int_{[a,b]}f\,\mathrm d\lambda_1$.
Nie korzystaliśmy tutaj z kryterium bezwzględnej
całkowalności KH, które udowodnimy poniżej.
\end{proof}
\begin{twierdzenie}[Kryterium bezwzględnej całkowalności KH]\label{ma1:kh:thm:kh-absolute}\label{it:thm:hk-absolute-integrability}
Niech $f:[a,b]\to\R$ będzie funkcją o skończonych wartościach na przedziale zwartym. Jeżeli $f\ge0$ jest KH-całkowalna, to $f\in\mathcal L^1$. Dla dowolnej takiej funkcji mamy
\[
f\in\mathcal L^1\quad\Longleftrightarrow\quad f\in\HK\ \text{i}\ |f|\in\HK.
\]
Na zwartym przedziale każda ograniczona funkcja jest KH-całkowalna wtedy i tylko wtedy, gdy jest mierzalna w sensie Lebesgue'a. W przypadku ograniczonej funkcji warunek konieczny i wystarczający całkowalności Riemanna to zerowa miara zbioru punktów nieciągłości (\atlasref{ME-08::thm:riemann-criterion}{ME-08: kryterium Lebesgue'a}).
\end{twierdzenie}
\begin{proof}
Dla nieujemnej $f\in\HK$ mierzalność wynika z twierdzenia \ref{ma1:kh:thm:kh-primitive-derivative}; obcięcia $h_n=\min(f,n)$ są mierzalne, ograniczone i całkowalne Lebesgue'a, więc należą także do KH. Z monotoniczności KH mamy $\int h_n\,\mathrm d\lambda_1\le(\HK)\int f<\infty$. Ponieważ $0\le h_n\uparrow f$ punktowo,
\atlasref{ME-08::thm:mct}{ME-08: twierdzenie Beppo Leviego}
daje $\int_{[a,b]}f\,\mathrm d\lambda_1=\lim_n\int_{[a,b]}h_n\,\mathrm d\lambda_1
\le(\HK)\int_a^b f<\infty$, czyli $f\in\mathcal L^1$. Zastosowanie pierwszej części do $|f|$ oraz inkluzja $\mathcal L^1\subset\HK$ dają równoważność. Równoważność dla ograniczonych funkcji KH wynika z mierzalności funkcji KH i skończoności miary przedziału. Dla $a<b$ kryterium całkowalności Riemanna jest dokładnie twierdzeniem \atlasref{ME-08::thm:riemann-criterion}{ME-08: kryterium Lebesgue'a}. Dla $a=b$ każda funkcja o skończonej wartości jest całkowalna Riemanna z całką zero, a zbiór punktów nieciągłości jest pusty.
\end{proof}
\begin{wniosek}[Części dodatnia i ujemna w kryterium KH]
\label{it:cor:hk-positive-negative-integrability}
Dla każdej funkcji rzeczywistej $f$ o skończonych
wartościach na $[a,b]$ następujące warunki są równoważne:
\begin{enumerate}[label=(\roman*)]
\item $f\in\mathcal L^1([a,b],\lambda_1)$;
\item $f^+,f^-\in\HK[a,b]$;
\item $f,|f|\in\HK[a,b]$;
\item $f\in\mathcal M[a,b]$.
\end{enumerate}
Jeżeli zachodzi którykolwiek z nich, to
\[
 (\HK)\int_a^b|f|=
 \int_{[a,b]}|f|\,\mathrm d\lambda_1
 =(\HK)\int_a^b f^++(\HK)\int_a^b f^-.
\]
Jeśli natomiast $f\in\HK[a,b]\setminus\mathcal L^1$,
to żadna z funkcji $f^+,f^-$ nie jest KH-całkowalna,
a obie ich całki Lebesgue'a wynoszą $+\infty$.
\end{wniosek}
\begin{proof}
Z (i) wynika (ii), ponieważ $0\le f^\pm\le|f|$,
więc obie części są Lebesgue'a całkowalne, a
$\mathcal L^1\subset\HK$ z twierdzenia
\ref{it:thm:mcshane-lebesgue} i jego dowodu.
Z (ii) wynika (iii) przez liniowość KH i tożsamości
$f=f^+-f^-$, $|f|=f^++f^-$. Implikacja
(iii)$\Rightarrow$(i) jest treścią twierdzenia
\ref{ma1:kh:thm:kh-absolute}, a równoważność
(i) i (iv) wynika z twierdzenia
\ref{it:thm:mcshane-lebesgue}. Wzór na całkę modułu
otrzymujemy z zgodności wartości całek i addytywności
całki nieujemnej.

Jeżeli $f\in\HK$ i jedna z funkcji $f^+,f^-$
należy do KH, to druga również należy do KH,
ponieważ $f^-=f^+-f$, $f^+=f+f^-$.
Z (ii)$\Rightarrow$(i) wynika wówczas $f\in\mathcal L^1$.
Dla $f\in\HK\setminus\mathcal L^1$ obie części są
mierzalne (twierdzenie \ref{ma1:kh:thm:kh-primitive-derivative}),
lecz żadna nie może mieć skończonej całki Lebesgue'a,
gdyż wówczas należałaby do KH. Stąd obie całki
Lebesgue'a są nieskończone. Całka Lebesgue'a ze
znakiem nie jest wtedy określona, mimo istnienia
skończonej całki KH.
\end{proof}

\begin{twierdzenie}[Kryterium oparte na funkcji pierwotnej]\label{ma1:kh:thm:ac-lebesgue-primitive}\label{it:thm:ac-lebesgue-primitive}
Funkcja $F:[a,b]\to\R$ jest bezwzględnie ciągła wtedy i tylko wtedy, gdy istnieje $f\in\mathcal L^1([a,b],\lambda_1)$ taka, że $F(x)=F(a)+\int_a^x f\,\mathrm d\lambda_1$ dla wszystkich $x$. Wówczas $F'=f$ prawie wszędzie. W szczególności, jeśli $g\in\HK[a,b]$ i $F(x)=F(a)+(\HK)\int_a^x g$ dla każdego $x\in[a,b]$, to
$g\in\mathcal L^1\Longleftrightarrow F\in\AC[a,b]$.
\end{twierdzenie}
\begin{proof}
Dla $a=b$ twierdzenie jest bezpośrednie: każda funkcja na singletonie jest absolutnie ciągła, całka jest zerowa, a równości prawie wszędzie są spełnione, ponieważ singleton ma miarę zero. Dalej zakładamy $a<b$. Jeśli $F(x)=F(a)+\int_{[a,x]} f\,\mathrm d\lambda_1$ z $f\in\mathcal L^1$, to dla skończonej rodziny przedziałów o rozłącznych wnętrzach zachodzi
$\sum_j|F(v_j)-F(u_j)|\le\int_{\bigcup_j[u_j,v_j]}|f|\,\mathrm d\lambda_1$. Ponieważ końce przedziałów mają miarę zero, miara ich sumy jest równa sumie długości. Lemat \ref{ma1:kh:lem:integral-ac} (wynik ME-17) dla tej sumy daje warunek bezwzględnej ciągłości $F$. Twierdzenie \atlasref{ME-25::thm:FTC-lebesgue-25}{ME-25: podstawowe twierdzenie rachunku całkowego} daje $F'=f$ prawie wszędzie.

Odwrotnie, niech $F$ będzie bezwzględnie ciągła. Wybieramy $\eta>0$ z definicji dla błędu $1$ i dzielimy $[a,b]$ na skończenie wiele odcinków o długościach mniejszych od $\eta$. Na każdym z nich dowolna suma modułów przyrostów wzdłuż podziału jest mniejsza od $1$. Jeżeli wybrano $m$ takich odcinków, to po wstawieniu ich końców do dowolnego podziału suma modułów przyrostów może jedynie wzrosnąć (nierówność trójkąta). Suma po zagęszczonym podziale jest mniejsza od $m$, ponieważ na każdym z $m$ odcinków jest mniejsza od $1$. Stąd wszystkie pierwotne sumy są ograniczone przez $m$. Zatem $V(x)=\Var_{[a,x]}F$ jest skończona. Addytywność wahania (lemat \ref{ma1:kh:lem:variation-additivity}) daje $V(v)-V(u)=\Var_{[u,v]}F$.

Wykażemy bezwzględną ciągłość $V$. Dla $\varepsilon>0$ wybieramy $\eta>0$ z warunku bezwzględnej ciągłości $F$ dla błędu $\varepsilon/2$. Jeżeli niepusta rodzina $[u_j,v_j]$, $1\le j\le m$, ma rozłączne wnętrza i sumę długości mniejszą od $\eta$, dla każdego $j$ wybieramy podział przybliżający wahanie na tym przedziale z błędem mniejszym od $\varepsilon/(2m)$. Wszystkie komórki tych podziałów mają rozłączne wnętrza i łączną długość mniejszą od $\eta$. Suma ich modułów przyrostów $F$ jest mniejsza od $\varepsilon/2$. Dodając błędy przybliżeń, otrzymujemy $\sum_j(V(v_j)-V(u_j))<\varepsilon$. To dowodzi $V\in\AC[a,b]$.

Funkcje $G_\pm=(V\pm(F-F(a)))/2$ są niemalejące, ponieważ $V(v)-V(u)\ge|F(v)-F(u)|$, i bezwzględnie ciągłe jako kombinacje liniowe funkcji bezwzględnie ciągłych. Przedłużamy je przez stałe poza $[a,b]$; przedłużenia są bezwzględnie ciągłe na każdym zwartym przedziale, gdyż komórki można rozciąć w $a,b$, a odcinki zewnętrzne nie wnoszą przyrostów. Obie funkcje wyznaczają na mocy \atlasref{ME-04::it:thm:stieltjes-construction}{ME-04: twierdzenia o konstrukcji miary Lebesgue'a--Stieltjesa} skończone dodatnie miary Lebesgue'a--Stieltjesa $\nu_\pm$ o przyrostach $\nu_\pm((u,v])=G_\pm(v)-G_\pm(u)$. Ponieważ $G_\pm(a)=0$ i przedłużenia są stałe
poza $[a,b]$, ich dystrybuanty są dokładnie równe
$G_\pm$: dla każdego $x$ z konstrukcji miar wynika
$\nu_\pm((-\infty,x])=G_\pm(x)$.
Z twierdzenia \atlasref{ME-25::thm:cdf-derivative-25}{ME-25: pochodna funkcji rozkładu} wynika
$\nu_\pm\ll\lambda_1$. Obie miary $\nu_\pm$ są skończone, a $\lambda_1$
jest $\sigma$-skończona. Stosujemy
\atlasref{ME-18::thm:RN-sigfin-18}{ME-18: twierdzenie Radona--Nikodyma}
na borelowskiej $\sigma$-algebrze prostej, otrzymując
nieujemne gęstości $g_\pm$ z
$\int_{\mathbb R}g_\pm\,\mathrm d\lambda_1=\nu_\pm(\mathbb R)<\infty$.
Wybieramy ich mierzalne reprezentanty o skończonych wartościach,
zmieniając ewentualne wartości nieskończone na zero
na mierzalnym zbiorze zerowym.
Dla $f=(g_+-g_-)|_{[a,b]}$ mamy
$\int_{[a,b]}|f|\,\mathrm d\lambda_1
\le\nu_+(\mathbb R)+\nu_-(\mathbb R)<\infty$.
Z definicji $G_\pm$ oraz reprezentacji ich miar,
dla każdego $x\in[a,b]$ otrzymujemy
\[
 F(x)-F(a)=G_+(x)-G_-(x)
 =\nu_+((a,x])-\nu_-((a,x])
 =\int_{[a,x]}f(t)\,\mathrm d\lambda_1(t).
\]
W ostatniej równości wykorzystujemy zerową miarę
singletonu $\{a\}$. Zatem $f\in\mathcal L^1([a,b],\lambda_1)$. Twierdzenie \atlasref{ME-25::thm:FTC-lebesgue-25}{ME-25: podstawowe twierdzenie rachunku całkowego} daje $F'=f$ prawie wszędzie.

Dla pierwotnej KH funkcji $g$, jeśli $g\in\mathcal L^1$, twierdzenie \ref{it:thm:mcshane-lebesgue} i inkluzja McShane'a w KH dają zgodność całek na każdym podprzedziale $[a,x]$. Z pierwszej części dowodu wynika bezwzględna ciągłość $F$. W przeciwnym kierunku reprezentacja Lebesgue'a daje $f=F'$ prawie wszędzie, a twierdzenie \ref{ma1:kh:thm:kh-primitive-derivative} o pochodnej pierwotnej KH daje również $g=F'$ prawie wszędzie. Zatem $g=f$ prawie wszędzie. Funkcja $g$ jest mierzalna i ma skończone wartości jako funkcja KH-całkowalna (twierdzenie \ref{ma1:kh:thm:kh-primitive-derivative}). Z niezmienniczości całki przy równości prawie wszędzie (\atlasref{ME-07::thm:ae-integral-invariance}{ME-07}) mamy $\int_{[a,b]}|g|\,\mathrm d\lambda_1=\int_{[a,b]}|f|\,\mathrm d\lambda_1<\infty$, czyli $g\in\mathcal L^1([a,b],\lambda_1)$.
\end{proof}
\subsection{Przestrzenie funkcji o ograniczonym wahaniu i absolutnie ciągłych}

\begin{twierdzenie}
\label{ma1:thm:bv-banach}
Dla $a<b$ przestrzeń $\BV[a,b]$ wszystkich rzeczywistych
funkcji o ograniczonym wahaniu (z wartościami w każdym
punkcie, bez utożsamiania prawie wszędzie) z normą
\[
 \|F\|_{\BV}=|F(a)|+\Var_{[a,b]}F
\]
jest przestrzenią Banacha i nie jest ośrodkowa.
Norma $\|F\|_\infty+\Var_{[a,b]}F$ jest jej równoważna.
\end{twierdzenie}
\begin{proof}
Liniowość przestrzeni i nierówność trójkąta dla wahania
wynikają z nierówności trójkąta dla każdej sumy przyrostów
i przejścia do supremum. Jednorodność jest bezpośrednia.
Jeśli norma znika, to $F(a)=0$ i
$|F(x)-F(a)|\le\Var_{[a,b]}F=0$ dla każdego $x$.
Ta sama nierówność daje $\|F\|_\infty\le\|F\|_{\BV}$,
skąd
\[
 \|F\|_{\BV}\le\|F\|_\infty+\Var_{[a,b]}F
 \le 2\|F\|_{\BV}.
\]
Niech $(F_n)$ będzie ciągiem Cauchy'ego w normie
$\|\cdot\|_{\BV}$. Jest jednostajnie Cauchy'ego,
więc punktowa zupełność $\R$ daje granicę $F$,
a przejście do granicy w oszacowaniu Cauchy'ego
daje zbieżność jednostajną.
Dla $\varepsilon>0$ wybieramy $N$ takie, że
$\|F_m-F_n\|_{\BV}\le\varepsilon$ dla $m,n\ge N$.
Ustalmy $n\ge N$ i dowolny podział $a=x_0<\cdots<x_j=b$.
Przechodząc $m\to\infty$ w skończonej sumie, otrzymujemy
\[
 |F(a)-F_n(a)|+
 \sum_{i=1}^{j}|(F-F_n)(x_i)-(F-F_n)(x_{i-1})|
 \le\varepsilon .
\]
Biorąc supremum po podziałach, otrzymujemy
$|F(a)-F_n(a)|+\Var_{[a,b]}(F-F_n)\le\varepsilon$.
W szczególności $F-F_n$ ma ograniczone wahanie.
Ponieważ $F_n\in\BV[a,b]$ i przestrzeń jest liniowa,
także $F\in\BV[a,b]$. Dopiero teraz możemy zapisać
$\|F-F_n\|_{\BV}\le\varepsilon$.
Z dowolności $\varepsilon$ wynika zupełność.
Dla $s\in(a,b)$ niech $H_s=\chi_{[s,b]}$.
Mamy $\Var_{[a,b]}H_s=1$. Dla $s<t$ różnica
$H_s-H_t$ przyjmuje kolejno wartości $0,1,0$,
a więc jej wahanie wynosi $2$ (podział zawierający
$a,s,t,b$ realizuje tę wartość; żaden podział nie
może mieć więcej niż dwa niezerowe przyrosty).
Nieprzeliczalna rodzina $(H_s)$ ma odległości
parami równe $2$. Rozłączne kule promienia $1/2$
nie mogą wszystkie przecinać przeliczalnego zbioru,
więc nie istnieje przeliczalny zbiór gęsty.
\end{proof}

\begin{twierdzenie}
\label{ma1:thm:ac-banach-separable}
Niech $a<b$. Przestrzeń $\AC[a,b]$ funkcji absolutnie ciągłych
z normą
\[
 \|F\|_{\AC}=|F(a)|+\int_{[a,b]}|F'|\,\mathrm d\lambda_1
\]
jest ośrodkową przestrzenią Banacha. Zachodzą
\[
 \AC[a,b]\cong\R\times L^1([a,b],\lambda_1),
 \qquad
 \Var_{[a,b]}F=\int_{[a,b]}|F'|\,\mathrm d\lambda_1
 \quad(F\in\AC[a,b]).
\]
Izomorfizm jest kanoniczny po ustaleniu lewego końca
przedziału i izometryczny dla normy sumy.
Ponadto $\AC[a,b]$ jest domkniętą właściwą
podprzestrzenią $\BV[a,b]$.
\end{twierdzenie}
\begin{proof}
Twierdzenie \ref{ma1:kh:thm:ac-lebesgue-primitive}
zapewnia istnienie pochodnej prawie wszędzie
oraz jej całkowalność. Symbol $[F']$ oznacza klasę
w $L^1$, po dowolnym mierzalnym uzupełnieniu
wartości na zbiorze zerowym, gdzie pochodna
nie istnieje. Liniowość odwzorowania
$F\mapsto(F(a),[F'])$ wynika z liniowości
różniczkowania poza sumą dwóch zbiorów zerowych.
Jego odwrotność ma postać
$(c,[f])\mapsto(x\mapsto c+\int_{[a,x]} f\,\mathrm d\lambda_1)$.
Wybór reprezentanta klasy $[f]$ nie wpływa na wartość
całki. Obie konstrukcje są wzajemnie odwrotne na mocy
reprezentacji funkcji absolutnie ciągłej i równości
pochodnej z funkcją podcałkową prawie wszędzie.
Przeniesienie normy sumy przez liniową bijekcję
daje normę $\|\cdot\|_{\AC}$ i izometrię.
Zupełność $L^1$ wynika z twierdzenia Riesza--Fischera
(\atlasref{ME-14::thm:riesz-fischer-14}{ME-14: twierdzenie Riesza--Fischera};
analiza funkcjonalna, AF 12, \atlasref{FA-12::thm:lp-banach12}{zupełność przestrzeni $L^p$}).
Dla ośrodkowości ustalmy $f\in L^1$ i $\varepsilon>0$.
Przez obcięcie wartości do $[-M,M]$ otrzymujemy błąd
w $L^1$ mniejszy od $\varepsilon/3$ dla dostatecznie
dużego $M$, z całkowalności $|f|$.
Zastępując wartości obciętej funkcji przez skończenie
wiele liczb wymiernych z błędem punktowym mniejszym
od $\varepsilon/(3(b-a))$, otrzymujemy funkcję prostą
$s=\sum_{i=1}^m c_i\chi_{A_i}$ z $c_i\in\Q$.
Z regularności wewnętrznej i zewnętrznej miary Lebesgue'a
(\atlasref{ME-04::thm:lebesgue-inner-regular}{ME-04: regularność wewnętrzna}
oraz \atlasref{ME-04::thm:lebesgue-outer-regular}{ME-04: regularność zewnętrzna}) wybieramy zbiory
zwarte $K_i\subseteq A_i$ i otwarte $O_i\supseteq A_i$
z $\lambda_1(O_i\setminus K_i)<\delta$.
Zwartość daje skończone pokrycie $K_i$ przedziałami
o wymiernych końcach zawartymi w $O_i$.
Ich sumę przecinamy z $[a,b]$, otrzymując $B_i$
z $\lambda_1(A_i\triangle B_i)<\delta$.
Mamy
\[
 \Big\|s-\sum_{i=1}^m c_i\chi_{B_i}\Big\|_{L^1}
 \le\delta\sum_{i=1}^m |c_i|.
\]
Wybór $\delta<\varepsilon/(3(1+\sum_i|c_i|))$
daje łączny błąd mniejszy od $\varepsilon$.
Rodzina tych skończonych sum jest przeliczalna
i składa się z funkcji schodkowych z dokładnością
do wartości w skończenie wielu punktach.
Iloczyn $\R\times L^1$ jest zatem zupełny i ośrodkowy,
gdyż norma sumy kontroluje obie współrzędne,
a iloczyn przeliczalnych zbiorów gęstych jest gęsty.

Niech $F(x)=F(a)+\int_{[a,x]}f\,\mathrm d\lambda_1$.
Dla dowolnego podziału suma modułów przyrostów
nie przekracza $\int_{[a,b]}|f|\,\mathrm d\lambda_1$.
Aby uzyskać odwrotną nierówność, wybieramy funkcję
schodkową $g$ stałą na wnętrzach komórek pewnego podziału
i taką, że $\|f-g\|_{L^1}<\varepsilon$; gęstość takich
funkcji wynika z poprzedniego argumentu.
Dla $G(x)=\int_{[a,x]}g\,\mathrm d\lambda_1$
suma modułów przyrostów na tym podziale jest równa
$\|g\|_{L^1}$. Z udowodnionej już górnej nierówności
$\Var_{[a,b]}(F-G)\le\|f-g\|_{L^1}$, więc
\[
 \Var_{[a,b]}F\ge\Var_{[a,b]}G-\Var_{[a,b]}(F-G)
 \ge\|g\|_{L^1}-\varepsilon
 \ge\|f\|_{L^1}-2\varepsilon.
\]
Dowolność $\varepsilon$ daje równość wahania z normą
pochodnej. W konsekwencji włożenie $\AC\subset\BV$
jest izometryczne. Zupełna podprzestrzeń przestrzeni
metrycznej jest domknięta: ciąg jej elementów zbieżny
w przestrzeni większej jest Cauchy'ego i ma granicę
w podprzestrzeni, a jednoznaczność granicy identyfikuje
obie granice. Właściwość wynika np. z obecności
nieciągłych funkcji $H_s$ w $\BV$.
\end{proof}

\begin{uwaga}[Związek z przestrzeniami Sobolewa]
Dla $1\le p\le\infty$ każda klasa $W^{1,p}(a,b)$
ma dokładnie jednego reprezentanta absolutnie ciągłego
na $[a,b]$ z pochodną w $L^p$; odwrotnie każda taka
funkcja wyznacza element $W^{1,p}(a,b)$.
Istnienie reprezentanta i zgodność pochodnych
wynikają z twierdzenia RRC 09
(\atlasref{PDE-09::thm:pde9-ac}{charakteryzacja przez bezwzględną ciągłość}).
Jednoznaczność na $[a,b]$ wynika z ciągłości:
dwie ciągłe funkcje równe prawie wszędzie we wnętrzu
muszą być tam równe w każdym punkcie, gdyż otwarte
otoczenie punktu niezgodności miałoby dodatnią miarę.
Przez ciągłość są równe także na końcach.
W szczególności dla $p=1$ otrzymujemy identyfikację
z $\AC[a,b]$, lecz z normami równoważnymi, a nie
identycznymi: pisząc $L=b-a$,
\[
 \|F\|_{L^1}\le L\bigl(|F(a)|+\|F'\|_{L^1}\bigr),
 \qquad
 |F(a)|\le L^{-1}\|F\|_{L^1}+\|F'\|_{L^1}.
\]
Pierwsze oszacowanie wynika z reprezentacji całkowej,
drugie z $|F(a)|\le |F(x)|+\|F'\|_{L^1}$ po
scałkowaniu po $x$ i podzieleniu przez $L$. Po dodaniu
$\|F'\|_{L^1}$ do obu nierówności dostajemy
\[
 \|F\|_{L^1}+\|F'\|_{L^1}\le(L+1)\|F\|_{\AC},
 \qquad \|F\|_{\AC}\le\max\{L^{-1},2\}
 (\|F\|_{L^1}+\|F'\|_{L^1}).
\]
To daje równoważność norm wprost.
W $\BV[a,b]$ rozważanej powyżej nie identyfikujemy
funkcji równych prawie wszędzie, dlatego nie należy
utożsamiać jej bez dodatkowych konwencji z przestrzenią
BV definiowaną za pomocą pochodnej dystrybucyjnej.
\end{uwaga}

\begin{przyklad}[Funkcja Cantora nie jest absolutnie ciągła]
Niech $K\subset[0,1]$ będzie klasycznym zbiorem
Cantora, a $\kappa$ bezatomową miarą probabilistyczną
skupioną na $K$, skonstruowaną w
\atlasref{ME-17::ex:cantor-measure-17}{ME-17: przykładzie miary Cantora}.
Jej dystrybuanta $C(x)=\kappa([0,x])$ jest niemalejąca
i ciągła (ciągłość miary z góry i z dołu oraz brak atomów).
Mamy $C(0)=0$ i $C(1)=1$. Poza $K$ funkcja $C$ jest
lokalnie stała, ponieważ każdy punkt spoza domkniętego
$K$ ma otoczenie o mierze $\kappa$ równej zero.
Ponieważ $\lambda_1(K)=0$, zachodzi $C'=0$ prawie wszędzie.
Nie może więc być absolutnie ciągła: w takim przypadku
twierdzenie \ref{ma1:kh:thm:ac-lebesgue-primitive}
dawałoby $1=\int_{[0,1]} C'\,\mathrm d\lambda_1=0$.
Co więcej, $C$ nie jest pierwotną KH żadnej funkcji
o skończonych wartościach: gdyby
$C(x)=C(0)+(\HK)\int_0^x f$, to z twierdzenia
\ref{ma1:kh:thm:kh-primitive-derivative} mielibyśmy
$f=C'=0$ prawie wszędzie. Lemat
\ref{ma1:kh:lem:kh-null-modification} zastosowany na
$[0,1]$ dawałby $C(1)=C(0)$, sprzecznie z konstrukcją.
\end{przyklad}
\subsection*{Konstrukcja całki McShane'a dla funkcji charakterystycznych}
\begin{twierdzenie}[Funkcje charakterystyczne zbiorów mierzalnych]\label{tw:mcshane-indykator}
Niech $E\subset[a,b]$ będzie zbiorem mierzalnym w sensie Lebesgue'a. Wówczas $\chi_E$ jest całkowalna w sensie McShane'a i
$(\mathcal M)\int_a^b\chi_E(t)\,\mathrm dt=\lambda_1(E)$.
\end{twierdzenie}
\begin{proof}
Ustalmy $\varepsilon>0$. Ponieważ $E$ jest mierzalny i $\lambda_1(E)\le b-a<\infty$, z
\atlasref{ME-04::thm:lebesgue-inner-regular}{ME-04: regularności wewnętrznej}
wybieramy zwarty $K\subset E$ z $\lambda_1(E\setminus K)<\varepsilon/2$,
a z \atlasref{ME-04::thm:lebesgue-outer-regular}{ME-04: regularności zewnętrznej}
otwarty zbiór $O'\supset E$ z $\lambda_1(O'\setminus E)<\varepsilon/2$.
Dla $O=O'\cap[a,b]$ mamy zatem $\lambda_1(O\setminus K)<\varepsilon$. Konstruujemy cechowanie: dla $t\in E$ wybieramy $\delta(t)>0$ tak, by $(t-\delta(t),t+\delta(t))\cap[a,b]\subset O$; dla $t\notin E$ dobieramy je tak, by wymienione otoczenie nie przecinało $K$ (jest to możliwe, ponieważ $K$ jest domknięty i $t\notin K$).

Niech $\mathcal P$ będzie dowolnym wolnym $\delta$-drobnym podziałem. Każda jego komórka z oznaczeniem $t\in E$ leży w $O$, a każda komórka ze znacznikiem $t\notin E$ nie przecina $K$. Ponieważ komórki pokrywają cały przedział, suma długości komórek ze znacznikami w $E$ mieści się między $\lambda_1(K)$ a $\lambda_1(O)$. Jest ona dokładnie sumą McShane'a $S(\chi_E,\mathcal P)$. Wobec $\lambda_1(K)\le\lambda_1(E)\le\lambda_1(O)$ otrzymujemy $|S-\lambda_1(E)|<\varepsilon$.
\end{proof}
\begin{wniosek}[Funkcje proste]\label{wn:mcshane-proste}
Niech $s=\sum_{k=1}^m c_k\chi_{E_k}$, gdzie $c_k\in\R$ i wszystkie $E_k\subset[a,b]$ są mierzalne. Wtedy $s\in\mathcal M[a,b]$ oraz $(\mathcal M)\int_a^b s(t)\,\mathrm dt=\sum_{k=1}^m c_k\lambda_1(E_k)$.
\end{wniosek}
\begin{proof}
Każdy składnik jest całkowalny na mocy twierdzenia \ref{tw:mcshane-indykator}; stosujemy liniowość całki McShane'a z twierdzenia \ref{it:thm:mcshane-basic} do skończonej kombinacji liniowej.
\end{proof}
\begin{przyklad}[Zbiór Cantora]
Dla klasycznego zbioru Cantora $K\subset[0,1]$ mamy $\lambda_1(K)=0$, więc $(\mathcal M)\int_0^1\chi_K(t)\,\mathrm dt=0$ na mocy twierdzenia \ref{tw:mcshane-indykator}. Ten sam wynik zachodzi dla funkcji Dirichleta z racji $\lambda_1(\Q\cap[0,1])=0$. Funkcja charakterystyczna klasycznego zbioru Cantora jest też całkowalna Riemanna, ponieważ poza domkniętym zbiorem Cantora jest lokalnie stała, więc zbiór jej nieciągłości ma miarę zero. Stosujemy kryterium Lebesgue'a (\atlasref{ME-08::thm:riemann-criterion}{ME-08: kryterium całkowalności Riemanna}). Funkcja Dirichleta nie jest całkowalna Riemanna, gdyż jest nieciągła w każdym punkcie.
\end{przyklad}
\subsection{Własności i twierdzenia graniczne}
\begin{wniosek}[Działania kratowe]\label{it:cor:mcshane-lattice}
Dla $f,g\in\mathcal M[a,b]$ także $f^+$, $f^-$,
$\max(f,g)$ i $\min(f,g)$ należą do $\mathcal M[a,b]$.
\end{wniosek}
\begin{proof}
Stosujemy twierdzenia \ref{it:thm:mcshane-basic}
i \ref{it:thm:mcshane-absolute} do tożsamości
$f^+=(|f|+f)/2$ i $f^-=(|f|-f)/2$
oraz $\max(f,g)=(f+g+|f-g|)/2$ i
$\min(f,g)=(f+g-|f-g|)/2$.
\end{proof}

\begin{twierdzenie}[Wspólne cechowanie]\label{it:thm:mcshane-equiintegrability}
Niech $f_n\in\mathcal M[a,b]$ z całkami $I_n$ i niech
$f_n(t)\to f(t)\in\R$ dla każdego $t$.
Załóżmy, że dla każdego $\varepsilon>0$ istnieje jedno
cechowanie $\delta$ takie, że $|S(f_n,\mathcal P)-I_n|<\varepsilon$
dla wszystkich $n$ i wszystkich wolnych $\delta$-drobnych
podziałów $\mathcal P$. Wówczas $f\in\mathcal M[a,b]$
oraz $I_n\to(\mathcal M)\int_a^b f$.
\end{twierdzenie}
\begin{proof}
Dla ustalonego $\varepsilon>0$ wybieramy wspólne cechowanie
dla błędu $\varepsilon/3$. Z lematu Cousina
\ref{ma1:kh:lem:kh-cousin} istnieje podział drobny
ze znacznikami należącymi do komórek, który jest
również wolnym podziałem drobnym.
Ustalmy taki podział. Jego skończona suma dla $f_n$ zbiega do sumy dla $f$,
ponieważ zbieżność zachodzi we wszystkich znacznikach.
Dla dostatecznie dużych $n,m$ moduł różnicy sum jest
mniejszy od $\varepsilon/3$, więc z nierówności trójkąta
$|I_n-I_m|<\varepsilon$. Niech $I=\lim_n I_n$.
Dla dowolnej dokładności $\eta>0$ wybieramy wspólne
cechowanie dla błędu $\eta/2$. Dla każdego wolnego
drobnego podziału $\mathcal Q$ przejście do granicy
w nierówności $|S(f_n,\mathcal Q)-I_n|<\eta/2$ daje
$|S(f,\mathcal Q)-I|\le\eta/2<\eta$.
\end{proof}

\begin{wniosek}[Zbieżność jednostajna]\label{it:cor:mcshane-uniform}
Jeżeli $f_n\in\mathcal M[a,b]$ i $f_n\to f$ jednostajnie,
to $f\in\mathcal M[a,b]$ i całki zbiegają.
\end{wniosek}
\begin{proof}
Dla dostatecznie dużych $n,m$ różnica $f_n-f_m$ jest
ograniczona. Dla dowolnego wolnego podziału mamy
$|S(f_n-f_m,\mathcal P)|\le(b-a)\|f_n-f_m\|_\infty$,
ponieważ długości komórek sumują się do $b-a$.
Z liniowości całki $f_n-f_m$ ma całkę $I_n-I_m$.
Dla dowolnego $\eta>0$ wybieramy cechowanie z definicji
tej całki dla błędu $\eta$ i drobny podział istniejący
na mocy lematu Cousina. Otrzymujemy
$|I_n-I_m|\le(b-a)\|f_n-f_m\|_\infty+\eta$.
Przechodząc z $\eta\downarrow0$, dostajemy
$|I_n-I_m|\le(b-a)\|f_n-f_m\|_\infty$.
Stąd $I=\lim I_n$ istnieje. Dla $\varepsilon>0$ wybieramy
$n$ tak, aby $(b-a)\|f-f_n\|_\infty<\varepsilon/3$
oraz $|I_n-I|<\varepsilon/3$. Na wolnych podziałach
drobnych względem cechowania dla $f_n$ z błędem
$\varepsilon/3$ mamy
\[
 |S(f,\mathcal P)-I|
 \le (b-a)\|f-f_n\|_\infty+|S(f_n,\mathcal P)-I_n|
     +|I_n-I|<\varepsilon.
\]
Nie zakładamy z góry całkowalności granicy ani ograniczoności
poszczególnych funkcji.
\end{proof}

\begin{wniosek}[Lokalność, zbieżność monotoniczna i zmajoryzowana]
\label{it:cor:mcshane-lebesgue-consequences}
Całkowalność McShane'a na $[a,b]$ pociąga całkowalność
na każdym podprzedziale; całka jest addytywna przy
sklejaniu dwóch sąsiednich przedziałów.
Jeśli $f_n\in\mathcal M[a,b]$ oraz $f_n\to f$ prawie
wszędzie, zbieżność całek i całkowalność $f$ zachodzą
w każdym z następujących przypadków:
\begin{enumerate}[label=(\alph*)]
\item $|f_n|\le g$ prawie wszędzie dla pewnej
$g\in\mathcal M[a,b]$;
\item $f_n\uparrow f$ prawie wszędzie i $\sup_n I_n<\infty$,
gdzie $I_n=(\mathcal M)\int_a^b f_n$.
\end{enumerate}
Zakładamy skończone wartości funkcji $f$; na zbiorze
wyjątkowym wolno nadać jej dowolne skończone wartości.
\end{wniosek}
\begin{proof}
Najpierw dowodzimy lokalności bez reprezentacji miarowej. Dla $a<c<b$ wybieramy cechowanie z kryterium Cauchy'ego na $[a,b]$. Dwa wolne drobne podziały $[a,c]$ uzupełniamy tym samym drobnym podziałem $[c,b]$, istniejącym dzięki lematowi Cousina \ref{ma1:kh:lem:kh-cousin}, zastosowanemu do ograniczenia cechowania do $[c,b]$. Po odjęciu wspólnej sumy kryterium Cauchy'ego daje całkowalność na $[a,c]$. Zamiana ról i powtórzenie argumentu daje całkowalność na każdym podprzedziale. Odwrotnie, jeżeli istnieją całki na obu częściach, wybieramy na nich cechowania dla błędu $\varepsilon/2$. Dla $t\ne c$ ograniczamy odpowiednie cechowanie przez $|t-c|/2$, a dla $t=c$ bierzemy minimum obu cechowań. Żadna komórka ze znacznikiem $t\ne c$ nie może leżeć po przeciwnej stronie $c$ niż jej znacznik, nawet gdy znacznik nie należy do komórki. Komórkę przecinającą $c$ rozcinamy w $c$, zachowując znacznik $c$. Po rozcięciu otrzymujemy wolne drobne podziały obu części o niezmienionej łącznej sumie. Nierówność trójkąta daje błąd mniejszy od $\varepsilon$ względem sumy całek lokalnych. To dowodzi również addytywności.

Dla twierdzeń granicznych stosujemy identyfikację \ref{it:thm:mcshane-lebesgue}: wszystkie funkcje McShane'a są całkowalne w sensie Lebesgue'a i odwrotnie. Wybieramy jeden mierzalny zbiór miary zero zawierający wszystkie zbiory wyjątkowe z założeń. Poza nim $f$ jest granicą funkcji mierzalnych; z zupełności miary Lebesgue'a wynika mierzalność $f$. W (a) stosujemy twierdzenie \atlasref{ME-08::thm:dct}{ME-08} z majorantą $|g|$. W (b) niech $N$ będzie wspólnym mierzalnym zbiorem zerowym,
poza którym zachodzą wszystkie nierówności i zbieżność.
Określamy $u_n=f_n-f_1$ i $u=f-f_1$ poza $N$,
a na $N$ nadajemy im wartość zero. Wówczas $0\le u_n\uparrow u$
wszędzie, a niezmienniczość całki przy zmianie na zbiorze zerowym
(\atlasref{ME-07::thm:ae-integral-invariance}{ME-07})
daje $\int_{[a,b]}u_n\,\mathrm d\lambda_1=I_n-I_1$.
Stąd $\sup_n\int_{[a,b]}u_n\,\mathrm d\lambda_1<\infty$.
Twierdzenie \atlasref{ME-08::thm:mct}{ME-08} daje całkowalność $u$ i zbieżność jego całek;
ponieważ $u=f-f_1$ prawie wszędzie, ten sam wniosek zachodzi dla
$f-f_1$. Dodając $f_1$, dostajemy całkowalność $f$ i zbieżność $I_n$.
Ponowna identyfikacja całek przez twierdzenie
\ref{it:thm:mcshane-lebesgue} kończy dowód.
\end{proof}
\endgroup
